\section{Related Work}
\label{sec:related}

In this section, we review prior work on knowledge-based robot planning, belief-space planning under uncertainty, human involvement in robot autonomy, and uncertainty-aware query systems.

\subsection{Knowledge-Based Robot Planning}
\label{sec:related_kb}

Knowledge-based robotic systems use structured representations to describe objects, properties, relations, actions, and task constraints. Logic-based knowledge systems and robot knowledge bases have been used to store observed facts, infer missing information, and support semantic reasoning in robot tasks~\cite{tenorth2012knowledge,liu2020graph}. For example, robots can infer likely object locations from qualitative spatial relations, container-location relations, or place-object co-occurrence knowledge~\cite{kunze2014using,hasegawa2023inferring}. Commonsense and relational reasoning further connect symbolic knowledge with perceptual evidence in partially observable environments~\cite{mota2021integrated,zhang2024icorpp,daruna2019robocse}. Ontologies and type hierarchies support consistent integration of new observations and generalization across related objects, actions, and task contexts~\cite{mitrevsk2021ontology,olivares2019review,manzoor2021ontology}.

Classical symbolic planning provides a computational layer on top of such knowledge representations. STRIPS and PDDL represent states as sets of symbolic facts and actions as preconditions and effects~\cite{fikes1971strips,aeronautiques1998pddl}. Related planning representations such as FDR and domain-independent planners expose compact state variables and enable efficient search over symbolic transition systems~\cite{helmert2009concise,hoffmann2001ff,helmert2006fast,correa2024lifted}. This structure is useful in robotics because high-level actions can be checked against symbolic preconditions and then grounded into executable robot skills. However, classical symbolic planning typically assumes that the current state is known, whereas real robot knowledge is often produced by noisy perception and partial observations.

Recent work has expanded robot knowledge beyond hand-coded symbolic rules using large language models and vision-language models. LLM-based planning systems use procedural and commonsense knowledge to propose task decompositions or action sequences, which are then grounded with robot affordances, programmatic interfaces, or execution feedback~\cite{ahn2022can,huang2022inner,singh2023progprompt,liang2023code,ding2023task}. Other systems construct richer environment knowledge from scene graphs and open-vocabulary spatial representations, allowing robots to reason about objects, places, and relations beyond a closed symbolic vocabulary~\cite{gu2024conceptgraphs,leanza2025conceptbot}. Retrieval-augmented methods use demonstrations, procedural knowledge, or task-relevant constraints to guide long-horizon robot planning~\cite{singh2025adaptbot,zhouenhancing,xu2024p,wang2025instructrag,kumar2026open}. Related approaches combine LLM-generated plans with symbolic constraints, spatial reasoning, or verification mechanisms to improve executability and consistency~\cite{huang2023voxposer,dai2024optimal,liu2025delta,paulius2025bootstrapping,quartey2025verifiably}.

These systems show that modern knowledge sources can greatly broaden what robots know and how they generate plans. Nevertheless, both traditional knowledge bases and foundation-model-based systems often commit inferred knowledge as a proposal, prior, or verified fact rather than maintaining an explicit execution-time belief over alternative symbolic worlds. Their outputs may remain uncertain because of perception errors, incomplete grounding, missing observations, or hallucinated semantic assumptions. Our work keeps the interpretability and action structure of symbolic planning while representing uncertain local facts as a belief over possible symbolic states. This lets the planner reason about what it does not know before committing perceptual outputs or inferred facts to the knowledge base.

\subsection{Belief-Space Planning and Symbolic Uncertainty}
\label{sec:related_uncertainty}

POMDPs provide a standard framework for decision making under partial observability~\cite{kaelbling1998planning}. Instead of assuming access to the true state, a POMDP maintains a belief distribution over possible states and selects actions to maximize expected return. Online methods such as POMCP scale this idea by using Monte Carlo tree search and particle beliefs over action-observation histories~\cite{silver2010monte}. Related online POMDP methods have been applied to robot navigation, object search, manipulation, exploration, and interaction under uncertainty~\cite{somani2013despot,lauri2016exploration,xiao2019objectsearch,pajarinen2022composition}.

These methods show that particle beliefs are effective when the robot must preserve multiple hypotheses about the world. However, many POMDP formulations operate over task-specific state variables rather than explicit symbolic knowledge bases. Conversely, symbolic planners expose action preconditions and logical structure, but often lack a principled belief representation. Prior work on symbolic planning under uncertainty addresses this gap by combining logical structure with probabilistic inference or POMDP-style planning~\cite{sanner2010symbolic,zhang2015corpp,serrano2021knowledge}. Such methods demonstrate that symbolic structure can help organize uncertainty, but uncertainty is usually resolved through autonomous sensing, transition models, or planning abstraction.

Our setting adds a different source of information: human semantic feedback. Some ambiguities cannot be reliably resolved by additional robot sensing alone, such as visually ambiguous object states or uncertain task-relevant categories. We therefore maintain a symbolic particle belief over reachable frontier states and allow the robot to query a human when the belief is insufficiently confident. The answer directly filters the planner's belief and updates the knowledge base.

\subsection{Human Involvement in Robot Autonomy}
\label{sec:related_human_loop}

Human involvement has long been used to improve robot reliability under uncertainty. In teleoperation and supervisory control, humans monitor the robot and intervene when needed~\cite{goodrich2008human,dragan2013legibility}. Human-on-the-loop systems preserve robot autonomy during nominal operation while relying on an operator to detect and correct failures~\cite{scharre2015introduction,nahavandi2017trusted,kelly2019hg,mandlekar2020human}. These systems can improve safety and robustness, but they require sustained human attention. The operator must notice that a problem exists, interpret the situation, and decide how to intervene, which can increase workload and reduce scalability in long-horizon or repetitive tasks~\cite{endsley2018automation,wong2017workload}.

Another approach is failure-triggered assistance, where the robot asks for help only after it fails, reaches a dead end, or detects an invalid state. This reduces continuous monitoring, but it treats human input as recovery rather than prevention. In tasks with irreversible or costly actions, asking after failure may be too late. A third approach is frequent or fixed querying, where the robot asks whenever a predefined uncertainty condition is met. This can improve task success, but it often produces unnecessary interactions because it does not distinguish uncertainty that matters for the current plan from uncertainty that is irrelevant.

Our work targets a proactive middle ground. The human is not asked to continuously supervise the robot, and the robot does not wait until failure. Instead, the robot estimates whether its current symbolic belief is confident enough to support planning. If not, it asks a targeted question about the symbolic fact expected to reduce the most uncertainty.

\subsection{Uncertainty-Aware Querying}
\label{sec:related_query}

Recent work has studied robots and embodied agents that ask for help based on uncertainty. Active learning and active preference learning query humans to improve models or infer preferences~\cite{sadigh2017active}. Proactive robot behavior has also been shown to improve collaboration when the robot can anticipate when assistance is useful~\cite{baraglia2017efficient}. More recent uncertainty-aware systems use model confidence, language-model uncertainty, or conformal prediction to decide when to request human input~\cite{knowno2023,liang2024introspective,mullen2025lbap}. These methods are important because they move beyond continuous supervision and show that selective querying can reduce unnecessary intervention.

However, many query systems focus on action-level uncertainty or model-output confidence. In such systems, the human may be asked which action the robot should take, whether a proposed action is acceptable, or whether a model prediction should be trusted. This differs from asking about the latent symbolic state that causes planning uncertainty. If the robot asks for an action, the human effectively participates in planning. If the robot asks for a state fact, the human provides information while the robot remains responsible for planning.

Our framework follows the second approach. The robot asks predicate-level questions about task-relevant symbolic facts, such as object properties, perceived categories, or local symbolic conditions. Human responses eliminate inconsistent hypotheses from the belief and update the knowledge base, after which planning proceeds autonomously. By integrating selective querying with symbolic belief maintenance and online planning, our approach treats human feedback as an information source for reducing planning uncertainty rather than as a mechanism for action selection.
